% Page 049

\seite{49}

\abschnitt{Ernte}
\pstart
die diesjährige Ernte kann als recht befriedigend angesehen wer\ohb
den, trotz der großen Trockenheit fiel der erste Gras- und Heu=\ob
schnitt wenigstens aus. Obst und Kartoffeln gaben es stellenweise\ob
wenig.

\pend
\abschnitt{Schulfeierlichkeiten}
\pstart

Der Geburtstag Sr. Majestät des Kaisers wurde von der Schuljugend\ob
gemeinschaftlich von den 4 Schulen der hiesigen Pfar=\ob
rei (Seffern) gefeiert, und die Feier hat sich schlicht gestaltet. Ab=\ob
wechselnd wurden Lieder von den größeren Kindern gesungen\ob
und Gedichte vorgetragen. Die Ansprache an die Kinder hielt\ob
Lehrer Heintzen aus Schleid. In einem zweiten Vortrage,\ob
gehalten von Lehrer Krämer aus Seffern, wurde das Hoch Kaiser\ob
Friedrichs ausgebracht.

\pend
\abschnitt{Osterprüfung}
\pstart

wurde am 28ten März durch den Ortsschulinspektor Herrn Pfarrer Zimmer aus

\pend
\begin{verse}
Seffern abgehalten.
\end{verse}
\pstart

\pend
\abschnitt{das Schuljahr 1912-13}
\pstart

Entlassen wurden in diesem Jahre 3 Kinder, 2 Knaben und\ob
1 Mädchen. Dieselben widmen sich der Landwirthschaft. Aufge=\ob
nommen wurden 5; 3 Mädchen und 2 Knaben. Die Schüler=\ob
zahl beträgt 23.\ Von diesen sind 10 Knaben und 13 Mädchen.\ob
Alle diese sind katholisch.

\margmark{Gesehen, 30. Mär\\Zimmer, Pfr.}%
z 1912.\ gesehen d.\ 5. 12.\ob
                       Lent

\pend
\abschnitt{Revision}
\pstart

Die hiesige Schule wurde am 6.\ Mai durch Herrn Schulrat\ob
Lentz aus Bitburg revidiert.
\pend
